\documentclass[11pt,a4paper]{article}
\usepackage{rankit}
\usepackage[margin=2.25cm,headheight=16pt]{geometry}
\usepackage{fancyhdr}
\usepackage{microtype}
\usepackage{caption}
\usepackage{enumitem}
\usepackage{lastpage}

\pagestyle{fancy}
\fancyhf{}
\fancyhead[L]{\RankItWordmark\quad\small RPI 2026/2027}
\fancyhead[R]{\small Organizacja projektu}
\fancyfoot[C]{\small\thepage}
\renewcommand{\headrulewidth}{0.4pt}
\renewcommand{\footrulewidth}{0pt}
\setlength{\parindent}{0pt}
\setlength{\parskip}{0.45em}
\setlist{nosep,leftmargin=*}
\newcolumntype{L}{>{\raggedright\arraybackslash}X}

\hypersetup{
  pdftitle={Organizacja i infrastruktura projektu RankIt},
  pdfauthor={Igor Jankowski, Piotr Ratkowski, Zachariasz Słupski, Tymon Jezionek}
}

\begin{document}

\begin{titlepage}
  \centering
  \vspace*{2.2cm}
  {\Huge\RankItWordmark\par}
  \vspace{0.7cm}
  {\Large Realizacja Projektu Informatycznego 2026/2027\par}
  \vspace{1.7cm}
  {\huge\bfseries\color{RankItNavy}Organizacja i infrastruktura projektu\par}
  \vspace{0.45cm}
  {\LARGE RankIt\par}
  \vspace{0.8cm}
  {\color{RankItTeal}\rule{0.72\textwidth}{1.5pt}\par}
  \vspace{1.3cm}
  \begin{tabular}{rl}
    \textbf{Identyfikator:} & RPI-RANKIT-ORG-001 \\
    \textbf{Wersja:} & 1.0 \\
    \textbf{Data:} & 24 października 2026 \\
    \textbf{Autorzy:} & Igor Jankowski, Piotr Ratkowski, \\
    & Zachariasz Słupski, Tymon Jezionek \\
  \end{tabular}
  \vfill
  {\large\textcolor{RankItNavy}{Repozytorium:} \href{https://github.com/elprofesoriqo/rpi}{RankIt na GitHubie}\par}
\end{titlepage}

\pagenumbering{Roman}
\tableofcontents
\clearpage
\pagenumbering{arabic}
\setcounter{page}{2}

\section*{Cel dokumentu}
\addcontentsline{toc}{section}{Cel dokumentu}
Dokument opisuje organizację projektu RankIt, sposób współpracy oraz miejsce przechowywania kodu i dokumentacji. Repozytorium GitHub jest wspólnym punktem dostępu do materiałów projektu. Określono podział odpowiedzialności członków zespołu oraz przyjęty sposób komunikacji.

\section{Opis projektu i produktu}
\subsection{Cel i problem}
RankIt to aplikacja webowa do tworzenia rankingów i oceniania różnych rzeczy. Ma zebrać propozycje i oceny w jednym miejscu oraz ułatwić porównywanie wyników. Projekt odpowiada na potrzebę prostego, wspólnego sposobu oceniania pomysłów, miejsc, produktów lub innych pozycji.

\subsection{Obszar zastosowania i odbiorcy}
Obszarem zastosowania jest grupowe porównywanie propozycji. Planowanymi odbiorcami są pojedynczy użytkownicy i małe grupy. To wstępne określenie rynku; zespół nie przeprowadził badań rynku ani nie przyjął modelu płatności.

\subsection{Proponowany zakres pierwszej wersji}
Użytkownik powinien móc utworzyć ranking, dodać do niego oceniane pozycje, przyznać im oceny i zobaczyć uporządkowane wyniki. Szczegóły, takie jak skala ocen, logowanie, współdzielenie rankingu i zasady edycji, zespół powinien ustalić przed implementacją.

Poza pierwszą wersją proponuje się pozostawić płatności, rekomendacje, powiadomienia i integracje z zewnętrznymi serwisami. Ogranicza to zakres do funkcji potrzebnych na demonstrację projektu.

\subsection{Ograniczenia, systemy i terminy}
Projekt jest realizowany w roku akademickim 2026/2027. Dokładny termin oddania nie został podany w opisie zadania. Nie wskazano systemów, z którymi aplikacja ma się integrować. Technologia aplikacji nie została jeszcze określona.

\begin{enumerate}
  \item Uzgodnienie wymagań i podziału obowiązków.
  \item Przygotowanie szkicu widoków i decyzja o modelu oceniania.
  \item Implementacja tworzenia rankingu i wystawiania ocen.
  \item Sprawdzenie scenariusza na przykładowych rankingach.
  \item Przygotowanie demonstracji i końcowych materiałów.
\end{enumerate}
Daty etapów należy dopasować do harmonogramu przedmiotu.

\section{Interesariusze i użytkownicy}
\begin{tabularx}{\textwidth}{p{0.29\textwidth}L}
  \toprule
  \textbf{Interesariusz} & \textbf{Rola i potrzeba} \\
  \midrule
  Zespół projektowy & Uzgadnia wymagania, buduje aplikację i przygotowuje dokumentację. Skład: Igor Jankowski, Piotr Ratkowski, Zachariasz Słupski i Tymon Jezionek. \\
  Opiekun przedmiotu & Ocenia zgodność projektu z wymaganiami RPI i udziela wskazówek. Nazwisko nie zostało przekazane. \\
  Twórcy rankingów & Przygotowują tematy i pozycje do oceny. \\
  Oceniający i odbiorcy & Wystawiają oceny lub korzystają z uporządkowanych wyników. \\
  \bottomrule
\end{tabularx}

\section{Zespół}
Skład zespołu: \textbf{Igor Jankowski, Piotr Ratkowski, Zachariasz Słupski i Tymon Jezionek}. Zrzut ekranu przekazany przez zespół pokazuje wszystkich czterech członków w grupie Messenger.

\begin{tabularx}{\textwidth}{p{0.28\textwidth}L}
  \toprule
  \textbf{Osoba} & \textbf{Obszar odpowiedzialności} \\
  \midrule
  Igor Jankowski & Koordynacja, utrzymanie repozytorium, dokumentacja i integracja zmian. \\
  Piotr Ratkowski & Projekt widoków i prace nad interfejsem użytkownika. \\
  Zachariasz Słupski & Logika rankingów i model danych. \\
  Tymon Jezionek & Plan testów, przypadki testowe i sprawdzanie funkcji aplikacji. \\
  \bottomrule
\end{tabularx}

Podział przypisuje każdej osobie konkretne zadania projektowe; nie stanowi oceny jej wcześniejszych kompetencji.

\begin{figure}[htbp]
  \centering
  \includegraphics[height=0.72\textheight,keepaspectratio]{../assets/rankit-team-messenger.png}
  \caption{Skład zespołu RankIt w grupie Messenger.}
\end{figure}

\section{Komunikacja w zespole i z interesariuszami}
Messenger jest kanałem bieżącej komunikacji zespołu. GitHub Issues służą do zapisywania zadań i decyzji, aby ustalenia dotyczące kodu pozostały przy repozytorium. Większe zmiany przechodzą przez Pull Request, a autor sprawdza podgląd różnic przed połączeniem.

Proponowany rytm pracy to krótkie, cotygodniowe spotkanie online trwające około 15 minut. Pytania do opiekuna zespół kieruje kanałem wskazanym przez prowadzącego przedmiot. Nie podano jeszcze nazwiska opiekuna ani terminów spotkań.

\section{Repozytorium kodu}
\begin{itemize}
  \item Adres: \href{https://github.com/elprofesoriqo/rpi}{Repozytorium RankIt na GitHubie}
  \item Osoba odpowiedzialna za repozytorium: Igor Jankowski.
  \item Widoczność: publiczna.
  \item Główna gałąź: \texttt{main}; zadania realizujemy na krótkich gałęziach \texttt{feature/...}, \texttt{fix/...} lub \texttt{docs/...}.
\end{itemize}

Repozytorium zawiera dokumenty organizacyjne. Kod aplikacji i konfiguracja budowania nie zostały jeszcze dodane, dlatego nie deklarujemy działającego CI ani zautomatyzowanych testów.

\section{Dokumentacja i szablon dokumentów}
Kod i dokumentację przechowujemy w jednym publicznym repozytorium GitHub. Katalog \texttt{docs/} zawiera źródła LaTeX w \texttt{latex/} oraz gotowe pliki PDF w \texttt{output/pdf/}. Obowiązuje schemat krótkich nazw z małych liter i łączników; wersję dokumentu zapisujemy w metryce, a historię zmian przechowuje Git.

Szablon projektu zawiera stronę tytułową, identyfikator, wersję, datę, autorów, historię zmian i miejsce na źródła. Stosuje granat \texttt{\#243B53} i turkus \texttt{\#2CB1BC}. Za porządek w dokumentacji odpowiada Igor Jankowski jako opiekun repozytorium.

\section{Narzędzia}
\begin{tabularx}{\textwidth}{p{0.31\textwidth}L}
  \toprule
  \textbf{Obszar} & \textbf{Narzędzie i zastosowanie} \\
  \midrule
  Kod i wersjonowanie & Git oraz GitHub. \\
  Zadania i przegląd zmian & GitHub Issues i Pull Requests. \\
  Komunikacja & Messenger; cotygodniowe spotkanie online jest proponowanym rytmem pracy. \\
  Dokumentacja & Źródła LaTeX i gotowe PDF-y w katalogu \texttt{docs/}. \\
  Tworzenie aplikacji & Stos technologiczny i narzędzia budowania wymagają decyzji zespołu. \\
  \bottomrule
\end{tabularx}

Na tym etapie nie dokładamy tablicy, CI ani dodatkowych usług, dopóki zespół nie ustali technologii i potrzeby.

\section{Wykorzystanie AI i wnioski}
AI wykorzystano do odczytania wymagań zadania, przygotowania układu raportu, zaproponowania prostych zasad repozytorium i redakcji pierwszej wersji dokumentów. Przykładowe polecenie z tej pracy: \emph{Przygotuj prostą dokumentację organizacji i infrastruktury projektu; aplikacja to RankIt do tworzenia rankingów i oceniania.} Wynikiem były: raport, wzór dokumentu, prezentacja oraz źródła LaTeX.

AI przyspieszyło przygotowanie struktury i tekstu, ale nie może potwierdzić umiejętności zespołu, terminów spotkań, nazwiska opiekuna ani wymagań aplikacji, których zespół nie podał. Członkowie projektu powinni sprawdzić propozycje przed oddaniem. Zespół dostarczył skład i rzeczywisty zrzut z komunikatora.

\subsection*{Informacje do potwierdzenia}
\begin{itemize}
  \item nazwisko opiekuna i termin oddania;
  \item podział obowiązków zgodny z umiejętnościami członków;
  \item skala ocen, logowanie i zasady udostępniania rankingów;
\end{itemize}

\section*{Podstawa opracowania}
\addcontentsline{toc}{section}{Podstawa opracowania}
Opis zadania \emph{Organizacja i infrastruktura projektu}, Realizacja projektu informatycznego 2026/2027.

\vfill
\begin{center}\small RankIt | RPI-RANKIT-ORG-001 | v1.0 | 24.10.2026\end{center}
\end{document}
